\chapter{Languages, automata and molecular recognition}
Does \texttt{1101} end in \texttt{01}? A human can inspect its final two bits.
A streaming machine does not know which bits will be last. It must update a
small memory after every symbol, then decide when input ends. Molecular
recognition adds another question: which physical configuration represents that
memory, and how do we observe it?

Chapter 15 generated words by rewriting. Here we recognize a language, then
extend the idea to a transducer that adds numbers. We assume finite sets,
probability vectors and the earlier strand-domain conventions.
The reference is an exact symbolic machine with a separately assigned noise
channel, not a molecular assay simulator.

\section{Remember only what a future suffix can distinguish}
\index{finite automaton}\index{recognition}
A deterministic finite automaton (DFA) is $A=(Q,\Sigma,\delta,q_0,F)$:
$Q$ is a finite state set, $\Sigma$ an alphabet,
$\delta:Q\times\Sigma\to Q$ a total transition function, $q_0$ an initial
state, and $F\subseteq Q$ the accepting states. For $w=a_1\cdots a_n$,
\begin{equation}
 q_{i+1}=\delta(q_i,a_{i+1}),\qquad
 w\in L(A)\ \Longleftrightarrow\ q_n\in F.
 \label{eq:dfa}
\end{equation}
One symbol causes one transition. Acceptance during a prefix is not necessarily
acceptance of the full word: input must finish.

For words ending in 01, use three states. State $q_0$ means no useful suffix;
$q_1$ means suffix 0; $q_2$ means suffix 01. A useful suffix is also a prefix
of the target pattern. The complete transition table is
\begin{center}
\begin{tabular}{ccl}\toprule
State & on 0 / on 1 & Meaning\\\midrule
$q_0$ & $q_1$ / $q_0$ & no useful suffix\\
$q_1$ & $q_1$ / $q_2$ & most recent bit is 0\\
$q_2$ & $q_1$ / $q_0$ & most recent two bits are 01\\\bottomrule
\end{tabular}
\end{center}
\Plate{01-dfa}{A complete recognizer for words ending in 01. Each edge consumes
its labeled symbol. The double circle marks acceptance at end-of-input.
There is exactly one outgoing transition per symbol from each state.}{fig:dfa}

The meanings hold by induction on prefix length. Appending 0 always gives
suffix 0. Appending 1 after suffix 0 completes 01. Appending 1 in either other
state leaves no useful suffix.

Three states are necessary, not just sufficient. Prefixes $\varepsilon$, 0
and 01 must be distinguishable: the empty continuation accepts only 01;
continuation 1 distinguishes 0 from $\varepsilon$. If two prefixes shared a
state, every future continuation would have the same acceptance outcome from
both, contradicting these witnesses. This is the practical core of
state minimization.

\section{Execute the rule before discussing chemistry}
\index{state invariant}\index{end-of-input}
\Listing{recognize}
The constant \texttt{DELTA} contains integer state indices in the table's order.
Validation rejects nonbinary characters; it does not silently strip whitespace,
because that would change the language. Empty input is legal and rejected by
this particular machine.
\begin{center}\input{\Generated/recognition.tex}\end{center}
\Plate{02-trace}{The cursor advances once per transition. Consuming 1,1,0,1
takes the initial state through 0,0,1,2. A state records a useful suffix,
not the number of processed symbols.}{fig:trace}

Tests enumerate every binary word through length nine and compare with the
independent string predicate \texttt{endswith("01")}. They also check each
prefix invariant, not only the final answer. A correct guess at the end of one
example cannot conceal an invalid intermediate transition.

Three running states need $\lceil\log_2 3\rceil=2$ conventional binary storage
bits. But the input has length $n$ and a full trace records $n+1$ states.
Constant working memory is not constant total storage when inputs and traces
are retained. Similarly, a molecular device can have few logical states while
input strands and transition reagents scale with the task.

\section{From recognition to a streaming transducer}
\index{transducer}\index{carry}
An automaton can emit output as well as change state. For binary addition,
read equal-width operands least-significant bit first. At position $i$,
$x_i,y_i,c_i\in\{0,1\}$ determine
\begin{align}
 t_i&=x_i+y_i+c_i,\\
 z_i&=t_i\bmod 2,\qquad c_{i+1}=\lfloor t_i/2\rfloor.
 \label{eq:carry}
\end{align}
The carry is the only persistent state. Output is now a sequence rather than
one membership bit; a transition label is $(x_i,y_i)/z_i$.
\Plate{03-transducer}{A carry transducer evaluates one column at a time.
Pairs 01 and 10 preserve the carry, while their output depends on that carry.
A final carry output is required after the last input column.}{fig:transducer}

After $k$ positions the arithmetic invariant is
\begin{equation}
 \sum_{i=0}^{k-1}2^i(x_i+y_i)
 =
 \sum_{i=0}^{k-1}2^i z_i+2^k c_k.
 \label{eq:addinv}
\end{equation}
At $k=0$ both sides vanish because $c_0=0$. For the induction step, substitute
$x_k+y_k+c_k=z_k+2c_{k+1}$, multiply by $2^k$, and cancel the old carry.
At width $n$, emitting $c_n$ as bit $n$ yields the exact sum. Dropping it yields
only the sum modulo $2^n$.
\Listing{add_bits}
For $13+3$, width-four inputs are \texttt{1011} and \texttt{1100}
in least-significant-first notation.
\begin{center}\input{\Generated/addition.tex}\end{center}
The output is \texttt{00001}, which reverses to conventional \texttt{10000},
or 16. Reading the input strings with conventional significance describes
different operands. The symbol alphabet and the number representation are
separate parts of the contract.

Exhaustive tests compare all operand pairs through width six against integer
addition. Zero-width input returns one carry bit 0. Unequal widths are rejected;
a caller may explicitly pad the shorter operand with high-order zero bits.

\section{Compile a state path into local molecular constraints}
\index{scaffold}\index{domain-level abstraction}
A logical transition need not correspond to a molecule stepping along a tape.
One can encode adjacent state consistency in an assembled structure.
At position $i$, a compute element carries an incoming state label, an
input-dependent transition and an outgoing state label. Neighboring elements
must agree on the shared state. Anchoring the initial state prevents a locally
consistent path from starting in an unintended state.
\Plate{04-scaffold}{Conceptual position-addressed molecular constraint scheme.
Compute strands bind antiparallel to scaffold addresses. Complementary state
domains enforce agreement across neighboring positions. The crossed alternative
marks an incompatible interface. This domain schematic is not a sequence design
or a prescribed kinetic pathway.}{fig:scaffold}

A forward DFA follows one path from $q_0$. An assembly interpretation may admit
many partial configurations while constraints resolve. Final state-path relations
can agree while kinetics, occupancy, resources and failures differ.
Local constraint satisfaction does not imply each molecule follows the sequential
DFA trace.

The September 2026 Scaffolded DNA Computer report presents an FSM-to-scaffold
compilation scheme and demonstrated arithmetic and nondeterministic programs
\cite{sdc}. Its main text distinguishes domain-level bind/replace models from
detailed pathways that are not experimentally enforced. This motivates comparing
semantics across physical architectures; it does not validate our figure's
domains or assigned noise parameters.

Audit initial-state enforcement, unique positions, domain complementarity,
end-of-input reporting and the meaning of each output. An accepting signal must
depend on the complete intended word. A missing final strand is not automatically
a logical reject state.

\section{Keep transition error separate from molecule loss}
\index{absorbing loss}\index{observation}
Let $p_i(q)$ be the probability of state $q$ after $i$ symbols. Add state
$\dagger$ for irrecoverable loss. At each step, loss occurs with assigned
probability $\ell$. Conditional on survival, the intended next state occurs
with probability $1-e$; each of the two other states receives $e/2$.
This symmetric error rule is a teaching assumption, not a fitted mechanism.

For symbol $a$, define a row-stochastic matrix
\begin{align}
 T_a(q,\dagger)&=\ell,\\
 T_a(q,r)&=(1-\ell)\left[(1-e)\mathbf1_{r=\delta(q,a)}
       +\frac e2\mathbf1_{r\ne\delta(q,a)}\right],\\
 T_a(\dagger,\dagger)&=1.
\end{align}
Other entries from $\dagger$ are zero. Then $p_{i+1}=p_iT_{a_{i+1}}$.
Every row sums to one, so total probability is conserved.
\Listing{noisy_distribution}

The loss probability after $n$ steps is
\begin{equation}
 p_n(\dagger)=1-(1-\ell)^n,
 \label{eq:loss}
\end{equation}
because surviving paths must survive every step. This is independent of
transition error in this specific model. Unconditional acceptance is $p_n(q_2)$;
acceptance conditional on survival is
$p_n(q_2)/(1-p_n(\dagger))$ when the denominator is nonzero.
These answer different questions.
\Plate{05-errors}{Computed unconditional accepting-state mass and absorbing
loss for accepted words $0^{n-1}1$, with assigned $e=0.02$, $\ell=0.01$.
These are model probabilities, not measured assay sensitivities.
Lost material is not counted as a logical reject state.}{fig:errors}

With zero error and loss, a point distribution follows the DFA exactly.
With $\ell=1$, every nonempty input is lost. Tests check these limits,
nonnegativity, normalization and Equation~\ref{eq:loss}.
A chemical channel could be sequence-dependent, correlated, reversible or
concentration-dependent. Those changes may invalidate the independent-loss formula.

Readout is still missing. To connect state populations to a detector signal $Y$,
we need $P(Y\mid q)$, calibration and a threshold. State acceptance probability
is not experimental accuracy without an input distribution and an observation
model. The observation lessons from Chapter 10 still apply.

\section{Finite-state power and its limits}
\index{regular language}
A finite automaton recognizes a regular language. Fixed finite memory cannot
remember an arbitrarily large unmatched count. For example,
$\{0^n1^n:n\geq0\}$ has no fixed finite DFA. Suppose a recognizer had $m$ states.
Among the $m+1$ prefixes $\varepsilon,0,\ldots,0^m$, two lengths $i<j$
reach the same state. Appending $1^i$ must give the same acceptance outcome
from both, but only $0^i1^i$ belongs to the language. Contradiction.

This does not prohibit finite-width experiments: a larger machine can encode
a specified maximum length. It prohibits using a bounded demonstration as
evidence of unbounded counting with fixed memory. Physical finite-state
machines also have finite input and reagent inventories; formal recognition
does not eliminate fabrication or readout cost.

\section{Exercises with worked answers}
\begin{enumerate}
\item \textbf{Predict.} Is 010 accepted? \emph{Answer:} No; the final state is
$q_1$, despite visiting $q_2$ earlier.
\item \textbf{Trace.} Give states for 001. \emph{Answer:} $q_0,q_1,q_1,q_2$.
\item \textbf{Prove.} Why cannot $q_0,q_1$ merge?
\emph{Answer:} Continuation 1 rejects from $q_0$ and accepts from $q_1$.
\item \textbf{Implement.} Recognize words containing 01 anywhere.
\emph{Answer:} Make $q_2$ absorbing for both symbols; test against substring
membership, including trailing 0.
\item \textbf{Count.} Contrast running-state and trace storage.
\emph{Answer:} Two bits suffice for the state, while a fixed-width trace of
$n+1$ states needs storage proportional to $n$.
\item \textbf{Derive.} Prove Equation~\ref{eq:addinv}.
\emph{Answer:} Use zero carry as the base case and cancel $2^kc_k$ after
substituting the local carry identity into the next partial sum.
\item \textbf{Overflow.} Add width-three 7 and 1.
\emph{Answer:} LSB-first 111 plus 100 yields 0001; discarding carry gives zero modulo eight.
\item \textbf{Boundary.} Interpret zero-width addition.
\emph{Answer:} Both operands are zero and only initial carry 0 is emitted.
\item \textbf{Probability.} Find survival after ten steps at $\ell=0.01$.
\emph{Answer:} $0.99^{10}\approx0.90438$; loss is approximately 0.09562.
\item \textbf{Counterexample.} Does zero accepting signal establish rejection?
\emph{Answer:} No; loss and detector failure can remove an accepting signal.
\item \textbf{Audit.} Test a state-domain compilation.
\emph{Rubric:} Check local transitions, initial constraint, ordering, final reporting,
mismatch controls and calibration. Symbolic parity does not validate sequences.
\item \textbf{Limit.} Why does success on $0^n1^n$ through $n=10$ not establish
all $n$? \emph{Answer:} A bounded counter handles those cases; repeated states
defeat a fixed DFA over an unbounded range.
\end{enumerate}
\section{Terms and reproducible boundary}
A \Term{recognizer} returns membership; a \Term{transducer} emits a sequence.
A \Term{state invariant} explains memory at each prefix.
A \Term{transition channel} assigns next-state probabilities.
The \Term{end-of-input} boundary is part of interpretation, not another data bit here.
The README reproduces code and tables. Later formal models ask which additional
memory changes computational power; later engineering chapters ask whether
molecular interfaces survive noise and composition.
